\documentclass[10pt,a4paper]{amsart}
\usepackage[margin=19mm]{geometry}
\usepackage{amsmath,amssymb,mathtools}
\usepackage[scr=rsfs,cal=euler]{mathalfa}
\usepackage{xfrac}
\usepackage[unicode,hidelinks,bookmarks=false]{hyperref}
\newcommand{\ie}{i.e.\ }
\newcommand{\cf}{cf.\ }
\newcommand{\resp}{respectively\ }
\newcommand{\Subsection}{\subsection}
\providecommand{\nobreakdash}{\nobreak\hbox{-}}
\setlength{\emergencystretch}{2em}
\multlinegap0pt
\usepackage{bm}
\usepackage{bbm}
\usepackage{dsfont}
\usepackage{xspace}
\usepackage{cancel}
\usepackage{mathtools}
\usepackage{amsthm}
\makeatletter
\@ifundefined{thm}{\newtheorem{thm}{Thm}}{}
\@ifundefined{lem}{\newtheorem{lem}[thm]{Lemma}}{}
\makeatother
\title[Resolution of singularities of the odd nilpotent cone of ort]{Resolution of singularities of the odd nilpotent cone of orthosymplectic Lie superalgebras}
\author{Ivan Motorin}
\date{Source: arXiv:2107.10007v4 (CC BY 4.0)}
\begin{document}
\maketitle

\noindent\emph{Source and licence.} Resolution of singularities of the odd nilpotent cone of orthosymplectic Lie superalgebras. Version arXiv:2107.10007v4 (CC BY 4.0), distributed under the Creative Commons Attribution 4.0 International licence, as recorded in the arXiv licence metadata for this exact version and shown on the abstract page https://arxiv.org/abs/2107.10007v4. Full rights notice: licenses/arxiv-2107.10007-CC-BY-4.0.txt.\par\smallskip

\noindent\textit{Editorial scope.} Counted unit: the Lem whose source label is \texttt{lem5} in the article (source0.tex lines 385--388, source label \texttt{lem5}), delivered together with the article's own complete original proof of it (source0.tex lines 390--410, one \texttt{proof} environment of 21 lines). No mathematics is written, repaired or re-derived: statement and proof are the article's own text line for line. Required context retained uncounted: none. Every definition, display and label of those lines is retained unchanged. Omitted from this package and not redistributed: 1--384, 389--389, 411--637. Nothing inside a retained excerpt is omitted or altered: every retained body is delivered exactly as the article's lines are written, so no omission receipt of any kind accompanies this package. Citations: the retained excerpts carry no citation command at all, so no bibliography entry of the article is transcribed and no reference is redistributed. Reference adaptation: none was needed and none was made; every reference inside the delivered unit resolves inside it, so no unresolved reference or citation is produced. Printed slips kept literal: no printed word, symbol, hypothesis or number of the counted statement or of its proof was altered. Layout and class: the article's own class is \texttt{amsart} with packages including cmap, inputenc, fontenc, amsmath, amssymb, amsthm, blkarray, graphicx, geometry, babel, euscript, tikz, tikz-cd, placeins; the wrapper is the lane's pinned \texttt{amsart} editorial wrapper at 10pt on A4, which restates the theorem-like environments the retained text needs and adds the article's own macro definitions that the retained text needs (none), with extra wrapper packages (bm, bbm, dsfont, xspace, cancel, mathtools). The delivered numbering is the unit's own. Assets: no retained span contains a figure environment, a graphics-inclusion command, a PDF-page inclusion, a table or a TikZ picture; no image file is copied into this package, the package declares no asset dependency and no image file is redistributed with it. Licence: version arXiv:2107.10007v4 of this article is distributed under the Creative Commons Attribution 4.0 International licence, as recorded in the arXiv licence metadata for this exact version and shown on the abstract page https://arxiv.org/abs/2107.10007v4. A full rights notice is delivered at licenses/arxiv-2107.10007-CC-BY-4.0.txt. Characters: every retained excerpt is pure ASCII, so the delivered engine, XeLaTeX with its default Latin Modern text font, reports no missing character.
\begin{lem}
  \label{lem5}
  The map $pr_1$ is birational.
\end{lem}

\begin{proof} For the birationality it is sufficient to show that for a regular nilpotent $A$ (with orbit of maximal dimension that is) there exists the unique flag such that $A$ respects it. Suppose $m=2n+1$, a regular nilpotent has the type of $\alpha_n$ (remember that we can directly compute dimensions of the orbit and the null cone)
\begin{equation}
    a_1\rightarrow b_1 \rightarrow \dots \rightarrow b_{2n} \rightarrow a_{2n+1} \rightarrow 0
\end{equation}
Since we have conditions $AF_0^{(1)}\subset F_1^{(0)}=\{0\}$, then $F_0^{(1)} = \langle a_{2n+1} \rangle$. Now $A^*F_1^{(1)}\subset F_0^{(1)}$ therefore $F_1^{(1)}$ has to be equal to $\langle b_{2n}\rangle$. Obviously we can continue this kind of reasoning and the flag is uniquely determined.\\
Let $m=2n+2$ then a regular nilpotent has the form $\alpha_n + \alpha_0$
\begin{equation}
    a_1\rightarrow b_1 \rightarrow \dots \rightarrow b_{2n} \rightarrow a_{2n+1} \rightarrow 0,  a_{2n+2} \rightarrow 0
\end{equation}
In this case $F_0^{(1)}$ lies in the kernel of $A$, however the first half of a flag must be self-orthogonal and $a_{2n+2}$ can not participate in it. This is why a flag is uniquely determined until the vectors $a_{n+2}$ and $b_{n+1}$. Now the only vectors that we can put in the middle component of a flag $F_0^{(n+1)}$ are either $a_{n+1}+ia_{2n+2}$ or $a_{n+1}+ia_{2n+2}$, besides the dual vector to already chosen one will be added in the $F_0^{n+2}$ component (we have strengthened the remark above as we only need to know the first half of the self-orthogonal flag and the connectedness component because of self-orthogonality condition). So this choice depends only on the connectedness component of the flag variety $F_0$.\\
Let $m=2n$. A regular nilpotent has the type $\beta_n + \alpha_0$
\begin{equation}
    b_1 \rightarrow a_1 \rightarrow \dots \rightarrow a_{2n-1} \rightarrow b_{2n} \rightarrow 0, a_{2n} \rightarrow 0
\end{equation}
Here $a$ and $b$ change roles since $AF_0^{(i)} \subset F_1^{(i)},A^*F_1^{(i)} \subset F_0^{(i-1)}(i\le n)$. That is why $F_1^{(1)}=\langle b_{2n}\rangle, F_0^{(1)}=\langle a_{2n-1}\rangle$ (again $a_{2n}$ does not participate because of self-orthogonality) and so on. Again, we can add $a_{n-1}+ia_{2n}$ or $a_{n-1}-ia_{2n}$ in the middle component $F_0^{(n)}$ and this corresponds to the choice of the connectedness component.\\
Finally, let $m=2n-1$. A regular nilpotent has the type $\beta_n$
\begin{equation}
    b_1 \rightarrow a_1 \rightarrow \dots \rightarrow a_{2n-1} \rightarrow b_{2n} \rightarrow 0,
\end{equation}
thus we see that $F_{1}^{1}= \langle b_{2n}\rangle, F_{0}^{1}=\langle a_{2n-1} \rangle$ and so on.
\end{proof}

\end{document}
